% =================================================================
% ==================== CONTINUATION: Q151--Q250 ===================
% =================================================================
\documentclass[aspectratio=169,8pt]{beamer}
\usetheme{Madrid}
\usepackage[utf8]{inputenc}
\usepackage[T1]{fontenc}
\usepackage{amsmath,amssymb}
\usepackage{booktabs}
\usepackage{enumitem}
\setlist[enumerate]{label=\Alph*.,leftmargin=*,itemsep=1pt,topsep=2pt}
\setlist[itemize]{leftmargin=*,itemsep=1pt,topsep=2pt}
\setbeamerfont{block body}{size=\tiny}
\setbeamerfont{block title}{size=\scriptsize}
\setbeamerfont{frametitle}{size=\footnotesize}
\setbeamerfont{normal text}{size=\tiny}
\setbeamertemplate{navigation symbols}{}
\setbeamertemplate{itemize item}{\tiny$\bullet$}
\setbeamertemplate{itemize subitem}{\tiny$\circ$}

\title[GARP RAI]{GARP Risk \& AI --- 250 Detailed Hard Questions}
\subtitle{Unofficial Exam Prep}
\author{Unofficial}
\date{\today}

\begin{document}
\section{Module 4 (cont.) and Module 5: Governance}

\begin{frame}{Q151 --- Medical Ethics Origins}
\begin{block}{Question}
Which document, adopted in 1947 in the aftermath of World War II, is considered foundational for modern medical ethics and, by extension, for the ethical treatment of human subjects in research that informs AI ethics today?

\begin{enumerate}
\item The Belmont Report.
\item The Nuremberg Code.
\item The Declaration of Helsinki.
\item The Declaration of Geneva.
\end{enumerate}  \end{block}
\begin{exampleblock}{Answer: B}\end{exampleblock}
\begin{block}{Explanation}
The Nuremberg Code (1947) established principles for ethical research involving human subjects, including voluntary consent. The Belmont Report came later (1978); the Declaration of Helsinki (1964) and Geneva (1948) built on the Nuremberg framework.
\end{block}
\end{frame}

\begin{frame}{Q152 --- Five AI Ethics Principles}
\begin{block}{Question}
Which of the following is NOT one of the five commonly cited AI ethics principles in the GARP Risk and AI curriculum?
\begin{enumerate}
\item Nonmaleficence.
\item Beneficence.
\item Profit maximisation.
\item Autonomy.
\end{enumerate}  \end{block}
\begin{exampleblock}{Answer: C}\end{exampleblock}
\begin{block}{Explanation}
The five principles are nonmaleficence, beneficence, justice, autonomy, and explainability. Profit maximisation is a business objective, not an ethical principle.
\end{block}
\end{frame}

\begin{frame}{Q153 --- Beneficence Limits}
\begin{block}{Question}
Which of the following best limits the ethical duty of beneficence?
\begin{enumerate}
\item There are no limits; beneficence is absolute.
\item Scarce resources, competing obligations, and the constraint that one cannot alleviate all suffering in the world.
\item Legal constraints only.
\item Financial constraints only.
\end{enumerate}  \end{block}
\begin{exampleblock}{Answer: B}\end{exampleblock}
\begin{block}{Explanation}
Beneficence is constrained by reasonableness: finite resources, competing moral demands, and the impossibility of universal assistance. Options A, C, and D are either too strong or partial.
\end{block}
\end{frame}

\begin{frame}{Q154 --- Autonomy and Surrogates}
\begin{block}{Question}
When a patient cannot make autonomous decisions (e.g., unconscious), who should make decisions on their behalf?
\begin{enumerate}
\item The AI system.
\item A surrogate decision-maker (family member, guardian) acting in the patient's best interests.
\item The hospital's CEO.
\item No one; decisions should be postponed indefinitely.
\end{enumerate}  \end{block}
\begin{exampleblock}{Answer: B}\end{exampleblock}
\begin{block}{Explanation}
When autonomy cannot be exercised directly, surrogates act in the patient's best interests. Options A, C, and D are not ethical or practical.
\end{block}
\end{frame}

\begin{frame}{Q155 --- GDPR Data Minimisation}
\begin{block}{Question}
A fintech company collects extensive personal data on its users because it might be useful someday. Which GDPR principle does this violate?
\begin{enumerate}
\item Data minimisation --- only data necessary for specified purposes should be collected.
\item Right to access.
\item Right to be forgotten.
\item Data portability.
\end{enumerate}  \end{block}
\begin{exampleblock}{Answer: A}\end{exampleblock}
\begin{block}{Explanation}
GDPR's data minimisation principle requires that personal data be adequate, relevant, and limited to what is necessary. Options B, C, and D are other GDPR rights unrelated to the collection issue.
\end{block}
\end{frame}

\begin{frame}{Q156 --- CCPA vs GDPR}
\begin{block}{Question}
Which of the following correctly distinguishes the CCPA from the GDPR?
\begin{enumerate}
\item CCPA applies only to California residents; GDPR applies to all EU residents and has extraterritorial reach.
\item CCPA covers all US residents.
\item GDPR is voluntary.
\item CCPA has no opt-out rights.
\end{enumerate}  \end{block}
\begin{exampleblock}{Answer: A}\end{exampleblock}
\begin{block}{Explanation}
CCPA is a California state law protecting California residents. GDPR applies across the EU (and to non-EU entities serving EU data subjects). Options B, C, and D are factually incorrect.
\end{block}
\end{frame}

\begin{frame}{Q157 --- AI Ethics in Financial Services}
\begin{block}{Question}
A bank's board asks the CRO to design an ethics framework for the bank's AI use in lending. Which of the following is the most appropriate first step?
\begin{enumerate}
\item Purchase an AI auditing tool.
\item Define the bank's AI ethics principles (fairness, transparency, accountability) and how they align with existing risk appetite and regulatory expectations.
\item Fire the data science team.
\item Adopt all AI techniques without restriction.
\end{enumerate}  \end{block}
\begin{exampleblock}{Answer: B}\end{exampleblock}
\begin{block}{Explanation}
Ethics frameworks start with principles aligned to business and regulatory context; tools come later. Options A, C, and D are premature or incorrect.
\end{block}
\end{frame}

\begin{frame}{Q158 --- Transparency vs Interpretability}
\begin{block}{Question}
Which statement best captures the difference between transparency and interpretability?
\begin{enumerate}
\item They are the same.
\item Transparency is about disclosure and openness (what the model does, what data it uses); interpretability is about how understandable the model's internal mechanisms are.
\item Transparency is only for regulators; interpretability is only for engineers.
\item Interpretability is a legal requirement.
\end{enumerate}  \end{block}
\begin{exampleblock}{Answer: B}\end{exampleblock}
\begin{block}{Explanation}
Transparency = openness about design, data, and decisions. Interpretability = degree to which internal logic is understandable. Options A, C, and D misstate the distinction.
\end{block}
\end{frame}

\begin{frame}{Q159 --- Governance Challenges: Power Asymmetries}
\begin{block}{Question}
Which of the following is a recognised governance challenge for AI ethics?
\begin{enumerate}
\item Power asymmetries between large tech firms and regulators.
\item Lack of interpretability.
\item Unpredictability of AI systems.
\item All of the above.
\end{enumerate}  \end{block}
\begin{exampleblock}{Answer: D}\end{exampleblock}
\begin{block}{Explanation}
All three are recognised governance challenges. Power asymmetries, opacity, and unpredictability complicate regulation and oversight.
\end{block}
\end{frame}

\begin{frame}{Q160 --- Lack of AI Ethics Structures}
\begin{block}{Question}
Compared with medical ethics, which structural element is generally missing in AI ethics?
\begin{enumerate}
\item Ethical codes.
\item Ethics committees.
\item Professional licensing and mandatory training.
\item All of the above are more mature in medicine than in AI.
\end{enumerate}  \end{block}
\begin{exampleblock}{Answer: D}\end{exampleblock}
\begin{block}{Explanation}
Medical ethics has mature codes, ethics committees, licensing, and mandatory training. AI ethics lacks comparable institutional structures. Options A, B, and C are each missing or less developed.
\end{block}
\end{frame}

\begin{frame}{Q161 --- Regulatory Landscape: EU AI Act}
\begin{block}{Question}
The EU AI Act came into force in 2024. Which statement about its scope is correct?
\begin{enumerate}
\item It regulates only medical AI.
\item It regulates AI systems based on a risk-tiered approach: unacceptable, high, limited, and minimal risk categories, with different obligations for each.
\item It bans all AI.
\item It applies only to EU-based entities.
\end{enumerate}  \end{block}
\begin{exampleblock}{Answer: B}\end{exampleblock}
\begin{block}{Explanation}
The AI Act uses a risk-tiered approach. Prohibited practices are banned; high-risk systems face strict obligations; limited-risk systems have transparency requirements; minimal-risk systems are mostly unregulated. Options A, C, and D misstate the Act.
\end{block}
\end{frame}

\begin{frame}{Q162 --- Regulatory Landscape: US}
\begin{block}{Question}
Which of the following best describes the current US federal regulatory landscape for AI?
\begin{enumerate}
\item A single comprehensive AI law.
\item No federal AI law; a patchwork of state laws, executive orders, and agency guidance, plus frameworks like the NIST AI RMF.
\item Only state-level AI regulation.
\item Complete prohibition of AI use in finance.
\end{enumerate}  \end{block}
\begin{exampleblock}{Answer: B}\end{exampleblock}
\begin{block}{Explanation}
As of the curriculum date, the US has no comprehensive federal AI law. Regulation is a patchwork of state laws (e.g., CCPA), executive orders, and agency frameworks (NIST AI RMF). Options A, C, and D are inaccurate.
\end{block}
\end{frame}

\begin{frame}{Q163 --- Regulatory Landscape: China}
\begin{block}{Question}
Which of the following best describes China's approach to AI regulation?
\begin{enumerate}
\item China has no AI regulation.
\item China has developed targeted regulations, including the Provisional Administrative Measures for Generative AI Services, the Personal Information Protection Law (PIPL), and algorithmic recommendation rules.
\item China only regulates medical AI.
\item China relies entirely on self-regulation by industry.
\end{enumerate}  \end{block}
\begin{exampleblock}{Answer: B}\end{exampleblock}
\begin{block}{Explanation}
China has adopted several targeted AI regulations, including the Generative AI Measures and PIPL. Options A, C, and D are factually incorrect.
\end{block}
\end{frame}

\begin{frame}{Q164 --- Relationship Between Ethics and Law}
\begin{block}{Question}
Which statement best captures the relationship between ethics and law in AI governance?
\begin{enumerate}
\item Ethics and law are identical.
\item Law establishes minimum standards of behaviour; ethics goes beyond law by guiding what is morally appropriate even when not legally required.
\item Ethics is subordinate to law.
\item Ethics replaces law.
\end{enumerate}  \end{block}
\begin{exampleblock}{Answer: B}\end{exampleblock}
\begin{block}{Explanation}
Ethics is broader than law: it reflects moral values and shapes laws but also guides behaviour where laws are silent. Options A, C, and D are incorrect.
\end{block}
\end{frame}



\begin{frame}{Q167 --- Data Quality}
\begin{block}{Question}
Which of the following is a critical concern when using alternative data for AI/ML models?
\begin{enumerate}
\item Alternative data is always more accurate than traditional data.
\item Alternative data may be unverified, subject to sudden definition changes, or may disappear, requiring continuous monitoring.
\item Alternative data never requires governance.
\item Alternative data is always free of bias.
\end{enumerate}  \end{block}
\begin{exampleblock}{Answer: B}\end{exampleblock}
\begin{block}{Explanation}
Alternative data quality can shift rapidly: providers may change definitions, cease data feeds, or introduce biases. Continuous monitoring is essential. Options A, C, and D are false.
\end{block}
\end{frame}


\begin{frame}{Q169 --- Data Classification}
\begin{block}{Question}
A firm classifies its data into Public, Internal, Confidential, and Restricted categories. What is the primary purpose of this classification?
\begin{enumerate}
\item To increase storage costs.
\item To determine handling, access controls, and protection levels appropriate to sensitivity and structure.
\item To confuse employees.
\item To comply with marketing requirements.
\end{enumerate}  \end{block}
\begin{exampleblock}{Answer: B}\end{exampleblock}
\begin{block}{Explanation}
Classification determines handling and protection based on sensitivity. Options A, C, and D are incorrect or nonsensical.
\end{block}
\end{frame}


\begin{frame}{Q171 --- Data Access Controls}
\begin{block}{Question}
Which access control model assigns permissions based on the user's role within the organisation?
\begin{enumerate}
\item Role-Based Access Control (RBAC).
\item Attribute-Based Access Control (ABAC).
\item Mandatory Access Control (MAC).
\item Discretionary Access Control (DAC).
\end{enumerate}  \end{block}
\begin{exampleblock}{Answer: A}\end{exampleblock}
\begin{block}{Explanation}
RBAC assigns permissions based on roles. ABAC uses attributes, MAC uses security labels, DAC lets owners set permissions. Option A is the correct answer.
\end{block}
\end{frame}




\begin{frame}{Q174 --- Model Governance Policies}
\begin{block}{Question}
Which of the following is typically included in model governance policies?
\begin{enumerate}
\item Definitions of a model and non-model.
\item Risk tiering methodology.
\item Documentation and validation requirements.
\item All of the above.
\end{enumerate}  \end{block}
\begin{exampleblock}{Answer: D}\end{exampleblock}
\begin{block}{Explanation}
Model governance policies define what constitutes a model, how to risk-rank it, and what documentation, validation, and monitoring are required. All three components are essential.
\end{block}
\end{frame}

\begin{frame}{Q175 --- Model Documentation}
\begin{block}{Question}
Which of the following is the primary purpose of model documentation in AI governance?
\begin{enumerate}
\item To provide a record of the model's design, assumptions, limitations, and usage, supporting validation, audit, and ongoing governance.
\item To increase model accuracy.
\item To reduce training time.
\item To eliminate the need for validation.
\end{enumerate}  \end{block}
\begin{exampleblock}{Answer: A}\end{exampleblock}
\begin{block}{Explanation}
Documentation supports validation, audit, reproducibility, and governance. It does not improve accuracy or eliminate validation. Options B, C, and D are incorrect.
\end{block}
\end{frame}

\begin{frame}{Q176 --- Model Inventory}
\begin{block}{Question}
Which of the following best describes the purpose of a model inventory?
\begin{enumerate}
\item To store the model's source code.
\item To maintain a registry of all models, their owners, risk tiers, validation status, and interdependencies, enabling enterprise-level oversight.
\item To train new models.
\item To visualise model performance.
\end{enumerate}  \end{block}
\begin{exampleblock}{Answer: B}\end{exampleblock}
\begin{block}{Explanation}
The model inventory is a registry for governance, tracking models across the organisation. Options A, C, and D are unrelated or partial.
\end{block}
\end{frame}

\begin{frame}{Q177 --- Model Landscape}
\begin{block}{Question}
What is a model landscape, and why is it important?
\begin{enumerate}
\item A landscape is a list of models. It is not important.
\item A landscape is a graphical or list-based representation of the interactions and interdependencies between models, revealing critical models and cascading risks.
\item A landscape is a data pipeline.
\item A landscape is a visualisation of model outputs.
\end{enumerate}  \end{block}
\begin{exampleblock}{Answer: B}\end{exampleblock}
\begin{block}{Explanation}
The model landscape maps interdependencies between models, revealing how failures can cascade. Options A, C, and D misstate the concept.
\end{block}
\end{frame}





\begin{frame}{Q180 --- Model Ownership}
\begin{block}{Question}
Who should typically own a model in an organisation's model governance framework?
\begin{enumerate}
\item The vendor who built the model.
\item A designated individual or committee in the business unit that uses the model and is accountable for its performance and use.
\item The IT department only.
\item The regulator.
\end{enumerate}  \end{block}
\begin{exampleblock}{Answer: B}\end{exampleblock}
\begin{block}{Explanation}
Model owners are typically in the business unit using the model, ensuring accountability. Options A, C, and D are not typical.
\end{block}
\end{frame}

\begin{frame}{Q181 --- Model Validation}
\begin{block}{Question}
When must model validation occur, and who should perform it?
\begin{enumerate}
\item Only after a model has caused a loss; performed by the developers.
\item Before initial deployment, periodically, and after material changes; performed by an independent validation function not involved in model development.
\item Only annually; performed by the same team that developed the model.
\item Never; validation is optional.
\end{enumerate}  \end{block}
\begin{exampleblock}{Answer: B}\end{exampleblock}
\begin{block}{Explanation}
Validation occurs pre-deployment, periodically, and after material changes, and must be independent. Options A, C, and D violate governance principles.
\end{block}
\end{frame}


\begin{frame}{Q184 --- Model Use Limits}
\begin{block}{Question}
Which of the following is the best practice for limiting model use?
\begin{enumerate}
\item Allow any use case without restriction.
\item Document approved uses and restrictions in the model inventory, and prevent use outside the validated domain (e.g., pricing one instrument with a model validated for another).
\item Restrict all models to a single use case regardless of validation.
\item Ignore model use restrictions.
\end{enumerate}  \end{block}
\begin{exampleblock}{Answer: B}\end{exampleblock}
\begin{block}{Explanation}
Model use limits should be documented and enforced to prevent misuse. Options A, C, and D violate governance best practice.
\end{block}
\end{frame}

\begin{frame}{Q185 --- Model Changes}
\begin{block}{Question}
When a model undergoes a material change (e.g., new data source, methodology update), what should happen in a well-governed framework?
\begin{enumerate}
\item The change is implemented immediately without review.
\item The change goes through a change-management process: impact assessment, documentation, validation, and approval.
\item The model is retired.
\item No action is needed.
\end{enumerate}  \end{block}
\begin{exampleblock}{Answer: B}\end{exampleblock}
\begin{block}{Explanation}
Material changes require impact assessment, documentation, validation, and approval. Options A, C, and D violate governance principles.
\end{block}
\end{frame}

\begin{frame}{Q186 --- Periodic Model Review}
\begin{block}{Question}
How frequently should AI/ML models be reviewed?
\begin{enumerate}
\item Never after initial deployment.
\item Only when performance degrades.
\item At least annually or more frequently if risk tier, usage, or environment warrants it; AI/ML models often require more frequent monitoring than traditional models.
\item Only every five years.
\end{enumerate}  \end{block}
\begin{exampleblock}{Answer: C}\end{exampleblock}
\begin{block}{Explanation}
Periodic review is required, with frequency based on risk tier. AI/ML models often warrant more frequent review due to complexity and data drift. Options A, B, and D are incorrect.
\end{block}
\end{frame}





\begin{frame}{Q189 --- Model Implementation Tasks}
\begin{block}{Question}
Which of the following is NOT a model implementation task?
\begin{enumerate}
\item Model design.
\item Core implementation.
\item System implementation.
\item Independent validation.
\end{enumerate}  \end{block}
\begin{exampleblock}{Answer: D}\end{exampleblock}
\begin{block}{Explanation}
Implementation tasks include model design, core implementation, and system implementation. Validation is separate. Options A, B, and C are implementation tasks.
\end{block}
\end{frame}

\begin{frame}{Q190 --- Model Adaptation}
\begin{block}{Question}
Core adaptation differs from system adaptation in that:
\begin{enumerate}
	\item Core adaptation modifies the model core (algorithm, parameters); system adaptation modifies the surrounding system (UI, data pipelines).
	\item They are the same.
	\item System adaptation is only for AI.
	\item Core adaptation is only for statistical models.
\end{enumerate}  \end{block}
\begin{exampleblock}{Answer: A}\end{exampleblock}
\begin{block}{Explanation}
	Core adaptation changes the model itself; system adaptation changes the supporting environment. Options B, C, and D misstate the distinction.
\end{block}
\end{frame}



\begin{frame}{Q192 --- Model Risk Reporting}
	\begin{block}{Question}
		What should model risk reporting to senior management and the board typically include?
		\begin{enumerate}
			\item Only the number of models in production.
			\item A summary of the model inventory, validation status, key risk issues, performance metrics, and remediation plans.
			\item The raw source code.
			\item Marketing material.
		\end{enumerate}  \end{block}
		\begin{exampleblock}{Answer: B}\end{exampleblock}
		\begin{block}{Explanation}
			Model risk reporting covers inventory, validation, risk issues, performance, and remediation. Options A, C, and D are incomplete or inappropriate.
		\end{block}
	\end{frame}
	
\begin{frame}{Q193 --- Model Change Management}
\begin{block}{Question}
Which of the following is a key component of model change management?
\begin{enumerate}
\item Informal verbal approval.
\item Structured change requests, impact assessment, documentation, validation, and approval.
\item Skipping validation for minor changes.
\item Never changing models.
\end{enumerate}  \end{block}
\begin{exampleblock}{Answer: B}\end{exampleblock}
\begin{block}{Explanation}
Change management includes structured processes: request, impact assessment, documentation, validation, approval. Options A, C, and D violate governance.
\end{block}
\end{frame}



\begin{frame}{Q195 --- GDPR Data Breach Notification}
\begin{block}{Question}
Under GDPR, how quickly must an organisation notify the supervisory authority of a personal data breach?
\begin{enumerate}
\item Immediately.
\item Within 72 hours of becoming aware of the breach, unless the breach is unlikely to result in a risk to rights and freedoms.
\item Within 30 days.
\item Only if the breach is publicly reported.
\end{enumerate}  \end{block}
\begin{exampleblock}{Answer: B}\end{exampleblock}
\begin{block}{Explanation}
GDPR Art. 33 requires notification within 72 hours. Options A, C, and D are inaccurate.
\end{block}
\end{frame}




\begin{frame}{Q198 --- Data Classification Levels}
\begin{block}{Question}
Which of the following is a typical data classification level?
\begin{enumerate}
\item Public, Internal, Confidential, Restricted.
\item Fast, Slow, Medium.
\item Red, Blue, Green only.
\item Alpha, Beta, Gamma.
\end{enumerate}  \end{block}
\begin{exampleblock}{Answer: A}\end{exampleblock}
\begin{block}{Explanation}
Common classification levels are public, internal, confidential, and restricted. Options B, C, and D are not standard classification levels.
\end{block}
\end{frame}

\begin{frame}{Q199 --- AI Ethics Framework}
\begin{block}{Question}
Which of the following is a key advantage of adopting an AI ethics framework?
\begin{enumerate}
\item It eliminates all AI risk.
\item It provides structured guidance for ethical decisions, enhances trust, and reduces regulatory non-compliance risk.
\item It maximises short-term profits.
\item It replaces the need for model validation.
\end{enumerate}  \end{block}
\begin{exampleblock}{Answer: B}\end{exampleblock}
\begin{block}{Explanation}
Frameworks provide guidance, build trust, and reduce regulatory risk. Options A, C, and D are incorrect.
\end{block}
\end{frame}


% =================================================================
\section{Additional Advanced Questions Q201--Q250}
% =================================================================

\begin{frame}{Q201 --- Inscrutability in Insurance}
\begin{block}{Question}
An insurance company uses a gradient-boosted model with hundreds of trees to price a complex commercial policy. The pricing actuary cannot fully explain how the model derives any particular premium. The risk function is concerned about ``inscrutability''. Which mitigation strategy is most aligned with the GARP Risk and AI curriculum?

\begin{enumerate}
\item Replace the model with a simple linear formula regardless of performance.
\item Keep the model but implement post-hoc explainability tools (SHAP, LIME), benchmark against simpler interpretable challenger models, and establish ongoing monitoring and documentation.
\item Deploy without any mitigation.
\item Assume inscrutability is acceptable.
\end{enumerate}  \end{block}
\begin{exampleblock}{Answer: B}\end{exampleblock}
\begin{block}{Explanation}
Inscrutability is best mitigated by explainability tools, benchmarking, and governance, rather than replacing models. Options A, C, and D are impractical or unsafe.
\end{block}
\end{frame}

\begin{frame}{Q202 --- Data Leakage in Time Series}
\begin{block}{Question}
A team building a forecasting model uses standard random k-fold cross-validation on 10 years of monthly economic data. The cross-validation performance is excellent, but the deployed model performs poorly. What is the most likely reason?
\begin{enumerate}
\item The model is underfit.
\item Standard k-fold shuffles data, allowing future information to leak into training. Time-series-aware validation (rolling/expanding window) should be used instead.
\item The data is corrupted.
\item The model is too simple.
\end{enumerate}  \end{block}
\begin{exampleblock}{Answer: B}\end{exampleblock}
\begin{block}{Explanation}
Time-series requires temporal validation to avoid leakage. Options A, C, and D are unlikely explanations.
\end{block}
\end{frame}

\begin{frame}{Q203 --- Model Risk: Model vs Non-Model}
\begin{block}{Question}
Which of the following is a qualified analytical tool (non-model) rather than a model under SR 11-7?
\begin{enumerate}
\item A Monte Carlo simulation for portfolio VaR.
\item A logistic regression for credit scoring.
\item A simple deterministic spreadsheet performing a fixed calculation.
\item A neural network.
\end{enumerate}  \end{block}
\begin{exampleblock}{Answer: C}\end{exampleblock}
\begin{block}{Explanation}
Non-models are tools with no material uncertainty or judgment (fixed calculations). Options A, B, and D are models with statistical uncertainty.
\end{block}
\end{frame}

\begin{frame}{Q204 --- Model Risk Ranking Factors}
\begin{block}{Question}
Which of the following is a standard factor in model risk ranking?
\begin{enumerate}
\item Materiality of model use.
\item Complexity of methodology.
\item Potential exposure (financial, regulatory, reputational).
\item All of the above.
\end{enumerate}  \end{block}
\begin{exampleblock}{Answer: D}\end{exampleblock}
\begin{block}{Explanation}
Model risk ranking considers materiality, complexity, and exposure. All three are standard factors.
\end{block}
\end{frame}



\begin{frame}{Q207 --- Model Validation: Outcomes Analysis}
\begin{block}{Question}
What is the purpose of outcomes analysis in model validation?
\begin{enumerate}
	\item To analyse model outputs against actual outcomes, checking that the model's predictions correspond to realised business or risk outcomes.
	\item To measure training time.
	\item To test the model's code line by line.
	\item To compare model speed across versions.
\end{enumerate}  \end{block}
\begin{exampleblock}{Answer: A}\end{exampleblock}
\begin{block}{Explanation}
	Outcomes analysis checks that model predictions align with reality. Options B, C, and D are unrelated.
\end{block}
															\end{frame}
															
										
																
								
																	
																	\begin{frame}{Q210 --- GenAI Model Monitoring}
																		\begin{block}{Question}
																			Why is continuous monitoring particularly important for GenAI systems?
																			\begin{enumerate}
																				\item Because GenAI models do not require monitoring.
																				\item Because outputs are non-deterministic, providers may update models, and drift can be subtle, requiring ongoing oversight.
																				\item Because GenAI systems are cheap to train.
																				\item Because GenAI systems require no validation.
																			\end{enumerate}  \end{block}
																			\begin{exampleblock}{Answer: B}\end{exampleblock}
																			\begin{block}{Explanation}
																				GenAI's non-determinism, vendor updates, and drift require continuous monitoring. Options A, C, and D are incorrect.
																			\end{block}
																		\end{frame}
																		
																	



\begin{frame}{Q213 --- Model Monitoring Metrics}
\begin{block}{Question}
Which of the following best describes the purpose of model monitoring?
\begin{enumerate}
\item To detect performance degradation, data drift, and concept drift, enabling timely intervention (retraining, recalibration, retirement).
\item To increase model complexity.
\item To eliminate the need for validation.
\item To market the model.
\end{enumerate}  \end{block}
\begin{exampleblock}{Answer: A}\end{exampleblock}
\begin{block}{Explanation}
Monitoring detects degradation and drift. Options B, C, and D are unrelated.
\end{block}
\end{frame}




\begin{frame}{Q217 --- Model Landscape Interdependencies}
\begin{block}{Question}
Why is identifying model interdependencies important in a model landscape?
\begin{enumerate}
\item To reduce model count.
\item Because one model's outputs may be another's inputs; a failure in one can cascade through others, creating systemic risk.
\item To speed up training.
\item To reduce costs.
\end{enumerate}  \end{block}
\begin{exampleblock}{Answer: B}\end{exampleblock}
\begin{block}{Explanation}
Interdependencies create cascading risks, which must be governed. Options A, C, and D are not the primary concern.
\end{block}
\end{frame}

\begin{frame}{Q218 --- Data Governance Committee}
\begin{block}{Question}
What is the typical role of a data governance committee?
\begin{enumerate}
\item Operational data entry.
\item Oversight of data strategy, policies, and practices, without direct operational responsibility.
\item Building models.
\item Marketing data products.
\end{enumerate}  \end{block}
\begin{exampleblock}{Answer: B}\end{exampleblock}
\begin{block}{Explanation}
The data governance committee provides oversight, not operational work. Options A, C, and D are operational roles.
\end{block}
\end{frame}

\begin{frame}{Q219 --- Explainability: Global vs Local}
\begin{block}{Question}
Which of the following best distinguishes global from local explainability?
\begin{enumerate}
	\item Global = overall model behaviour; local = specific prediction.
	\item Global = specific prediction; local = overall behaviour.
	\item They are identical.
	\item Local is for tree models only.
\end{enumerate}  \end{block}
\begin{exampleblock}{Answer: A}\end{exampleblock}
\begin{block}{Explanation}
	Global = overall; local = per-prediction. Options B, C, and D misstate the distinction.
\end{block}
\end{frame}

\begin{frame}{Q220 --- Model Risk: External Challenges}
\begin{block}{Question}
	Which of the following is NOT typically a source of model risk?
	\begin{enumerate}
		\item Poor data quality.
		\item Inappropriate model assumptions.
		\item Implementation errors.
		\item Increasing model training time.
	\end{enumerate}  \end{block}
	\begin{exampleblock}{Answer: D}\end{exampleblock}
	\begin{block}{Explanation}
		Training time is a performance concern, not a source of model risk. Poor data, wrong assumptions, and implementation errors are all recognised sources.
	\end{block}
\end{frame}

	
	\begin{frame}{Q222 --- Data Classification and Handling}
		\begin{block}{Question}
			Which of the following best describes the purpose of data classification in a governance framework?
			\begin{enumerate}
				\item To categorise data by sensitivity and structure so that appropriate handling, access, and protection controls can be applied.
				\item To increase data storage costs.
				\item To reduce data volume.
				\item To remove data.
			\end{enumerate}  \end{block}
			\begin{exampleblock}{Answer: A}\end{exampleblock}
			\begin{block}{Explanation}
				Classification drives controls. Options B, C, and D are incorrect.
			\end{block}
		\end{frame}
		
		\begin{frame}{Q223 --- Model Risk: Bias in AI}
			\begin{block}{Question}
				Which of the following is NOT a recognised source of algorithmic bias?
				\begin{enumerate}
					\item Problem specification and feature selection.
					\item Data collection and composition.
					\item Model development choices.
					\item Increasing the number of training epochs.
				\end{enumerate}  \end{block}
				\begin{exampleblock}{Answer: D}\end{exampleblock}
				\begin{block}{Explanation}
					Training epochs affect convergence, not fundamental bias sources. Sources of bias are problem specification, data, model development, and deployment.
				\end{block}
			\end{frame}
			
			\begin{frame}{Q224 --- Model Risk: Fairness Trade-offs}
				\begin{block}{Question}
					When base rates differ between groups, can all fairness metrics (demographic parity, predictive rate parity, equal opportunity) be satisfied simultaneously?
					\begin{enumerate}
						\item Yes, always.
						\item No, they are mathematically incompatible under unequal base rates.
						\item Only for logistic regression.
						\item Only for tree models.
					\end{enumerate}  \end{block}
					\begin{exampleblock}{Answer: B}\end{exampleblock}
					\begin{block}{Explanation}
						The impossibility theorem (Kleinberg, Chouldechova) shows these cannot all hold under unequal base rates. Options A, C, and D are incorrect.
					\end{block}
				\end{frame}
			
					
				
																																		
																																		\begin{frame}{Q227 --- GenAI Hallucinations}
																																			\begin{block}{Question}
																																				Which of the following is the most effective mitigation for LLM hallucinations in a high-stakes risk-management application?
																																				\begin{enumerate}
																																					\item Increase temperature.
																																					\item Use retrieval-augmented generation (RAG) to ground outputs in authoritative sources, and add human verification for critical decisions.
																																					\item Use a smaller model.
																																					\item Ignore hallucinations.
																																				\end{enumerate}  \end{block}
																																				\begin{exampleblock}{Answer: B}\end{exampleblock}
																																				\begin{block}{Explanation}
																																					RAG grounds model responses in retrieved content; human verification catches remaining errors. Options A, C, and D are inadequate or incorrect.
																																				\end{block}
																																			\end{frame}
																																			
																																				
																															
	
																																						
																																						\begin{frame}{Q231 --- AI Risk: Interconnectedness}
																																							\begin{block}{Question}
																																								An organisation has 50 models; outputs of 10 feed into inputs of 20 others. Which governance artefact is designed to capture these interactions?
																																								\begin{enumerate}
																																									\item Model inventory.
																																									\item Model landscape.
																																									\item Data dictionary.
																																									\item Validation report.
																																								\end{enumerate}  \end{block}
																																								\begin{exampleblock}{Answer: B}\end{exampleblock}
																																								\begin{block}{Explanation}
																																									The model landscape captures interdependencies. The inventory lists models; the data dictionary describes data; the validation report documents validation.
																																								\end{block}
																																							\end{frame}
																																					
																																								
																																										
																																									\begin{frame}{Q234 --- Model Risk: Model Use Limits}
																																										\begin{block}{Question}
																																											Why is it important to enforce limits on model use?
																																											\begin{enumerate}
																																												\item Because all uses are equivalent.
																																												\item Because using a model outside its validated scope (e.g., different products, geographies, or time horizons) can produce materially incorrect results and unexpected risks.
																																												\item Because regulators forbid all models.
																																												\item Because model use has no implications.
																																											\end{enumerate}  \end{block}
																																											\begin{exampleblock}{Answer: B}\end{exampleblock}
																																											\begin{block}{Explanation}
																																												Use limits prevent models from being applied outside their validated domain. Options A, C, and D are incorrect.
																																											\end{block}
																																										\end{frame}
																																										
							
																																											\begin{frame}{Q236 --- Model Risk: GenAI Governance Levers}
																																												\begin{block}{Question}
																																													Which of the following is an internal governance lever for GenAI?
																																													\begin{enumerate}
																																														\item Assigning clear ownership and accountability across the GenAI lifecycle.
																																														\item Tracking changes in the external regulatory environment.
																																														\item Relying on external standards only.
																																														\item Outsourcing all risk management.
																																													\end{enumerate}  \end{block}
																																													\begin{exampleblock}{Answer: A}\end{exampleblock}
																																													\begin{block}{Explanation}
																																														Internal levers include ownership, review processes, and escalation paths. Option B is external. Options C and D are not governance levers.
																																													\end{block}
																																												\end{frame}
																																												
																																												\begin{frame}{Q237 --- Model Risk: GenAI External Constraints}
																																													\begin{block}{Question}
																																														Which of the following is an external governance constraint for GenAI?
																																														\begin{enumerate}
																																															\item Evolving regulatory requirements (e.g., EU AI Act).
																																															\item Contractual terms with model providers.
																																															\item Emerging standards (e.g., NIST AI RMF, ISO/IEC 42001).
																																															\item All of the above.
																																														\end{enumerate}  \end{block}
																																														\begin{exampleblock}{Answer: D}\end{exampleblock}
																																														\begin{block}{Explanation}
																																															All three are external governance constraints. Regulatory, contractual, and standards-based constraints shape GenAI governance.
																																														\end{block}
																																													\end{frame}
																																													
																																																			
																									\begin{frame}{Q239 --- Model Risk: Ethical AI Frameworks}
																										\begin{block}{Question}
																											Which of the following is a practical benefit of an ethical AI framework in a financial institution?
																											\begin{enumerate}
																												\item Improved ability to evaluate ethical dilemmas, increased trust, and reduced regulatory risk.
																												\item Increased model accuracy.
																												\item Reduced training time.
																												\item Elimination of all AI risks.
																											\end{enumerate}  \end{block}
																											\begin{exampleblock}{Answer: A}\end{exampleblock}
																											\begin{block}{Explanation}
																												Ethics frameworks help evaluate dilemmas, build trust, and reduce regulatory risk. Options B, C, and D are not the primary benefits.
																											\end{block}
																										\end{frame}
																										
																										\begin{frame}{Q240 --- Model Risk: Explainability Techniques}
																											\begin{block}{Question}
																												Which of the following is a commonly used post-hoc explainability technique for complex models?
																												\begin{enumerate}
																													\item Linear regression.
																													\item SHAP values.
																													\item Logistic regression.
																													\item Ridge regression.
																												\end{enumerate}  \end{block}
																												\begin{exampleblock}{Answer: B}\end{exampleblock}
																												\begin{block}{Explanation}
																													SHAP is a post-hoc technique based on Shapley values. Options A, C, and D are modelling techniques, not explainability methods.
																												\end{block}
																											\end{frame}
																											
																											\begin{frame}{Q241 --- Model Risk: Data Classification in Practice}
																												\begin{block}{Question}
																													A firm stores customer Social Security Numbers alongside publicly available firm financials. How should these be classified?
																													\begin{enumerate}
																														\item Both public.
																														\item SSNs: restricted/confidential; financials: public.
																														\item Both restricted.
																														\item Both internal.
																													\end{enumerate}  \end{block}
																													\begin{exampleblock}{Answer: B}\end{exampleblock}
																													\begin{block}{Explanation}
																														SSNs are highly sensitive (restricted/confidential). Public financials are public. Options A, C, and D misclassify.
																													\end{block}
																												\end{frame}
																												
																												\begin{frame}{Q242 --- Model Risk: Metadata in Practice}
																													\begin{block}{Question}
																														A data scientist needs to know the exact definition of a column ``transactionamt'' in a large database. Where is this information most likely stored?
																														\begin{enumerate}
																															\item In the raw data.
																															\item In metadata (data dictionary).
																															\item In the model.
																															\item In the audit log.
																														\end{enumerate}  \end{block}
																														\begin{exampleblock}{Answer: B}\end{exampleblock}
																														\begin{block}{Explanation}
																															Metadata (data dictionary) describes column definitions, types, and relationships. Options A, C, and D are unrelated.
																														\end{block}
																													\end{frame}
																													
																													\begin{frame}{Q243 --- Model Risk: Access Controls in Practice}
																														\begin{block}{Question}
																															Which of the following is an example of a role-based access control decision?
																															\begin{enumerate}
																																\item All employees can access all data.
																																\item Only members of the credit risk team can access the credit risk database, based on their role.
																																\item All employees can access personal data if they sign a form.
																																\item Access is decided by seniority only.
																															\end{enumerate}  \end{block}
																															\begin{exampleblock}{Answer: B}\end{exampleblock}
																															\begin{block}{Explanation}
																																RBAC grants access based on roles. Options A, C, and D do not reflect RBAC principles.
																															\end{block}
																														\end{frame}
																														
																														\begin{frame}{Q244 --- Model Risk: Validation Activities}
																															\begin{block}{Question}
																																Which of the following is a typical model validation activity?
																																\begin{enumerate}
																																	\item Assessment of conceptual soundness.
																																	\item Review of data appropriateness.
																																	\item Review of implementation and outcomes analysis.
																																	\item All of the above.
																																\end{enumerate}  \end{block}
																																\begin{exampleblock}{Answer: D}\end{exampleblock}
																																\begin{block}{Explanation}
																																	Validation covers conceptual soundness, data, implementation, and outcomes. All three are standard activities.
																																\end{block}
																															\end{frame}
																															
																															\begin{frame}{Q245 --- Model Risk: AI/ML Model Ranking}
																																\begin{block}{Question}
																																	Which of the following would most likely elevate an AI/ML model to a higher risk tier?
																																	\begin{enumerate}
																																		\item High complexity and use in a regulatory capital calculation.
																																		\item Frequent use across business units with material financial impact.
																																		\item Use with alternative data sources introducing fairness and privacy concerns.
																																		\item All of the above.
																																	\end{enumerate}  \end{block}
																																	\begin{exampleblock}{Answer: D}\end{exampleblock}
																																	\begin{block}{Explanation}
																																		Complexity, materiality, and novel data sources all increase risk tier. All three elevate the ranking.
																																	\end{block}
																																\end{frame}
																																
																																\begin{frame}{Q246 --- Model Risk: Model Risk Reporting}
																																	\begin{block}{Question}
																																		Which of the following should a model risk report to the board typically include?
																																		\begin{enumerate}
																																			\item The full model inventory, validation status, key risk issues, and remediation plans.
																																			\item The full source code of all models.
																																			\item Detailed marketing plans.
																																			\item Employee salaries.
																																		\end{enumerate}  \end{block}
																																		\begin{exampleblock}{Answer: A}\end{exampleblock}
																																		\begin{block}{Explanation}
																																			Reports include inventory, validation status, risk issues, and remediation. Options B, C, and D are inappropriate.
																																		\end{block}
																																	\end{frame}
																																	
																																	\begin{frame}{Q247 --- Model Risk: Data Minimisation}
																																		\begin{block}{Question}
																																			A marketing team wants to collect more customer data than is strictly necessary for a new feature. What does the data minimisation principle advise?
																																			\begin{enumerate}
																																				\item Collect as much data as possible to maximise future options.
																																				\item Collect only what is necessary for the specified purpose; unnecessary data should not be collected.
																																				\item Collect data without any governance.
																																				\item Ignore the principle.
																																			\end{enumerate}  \end{block}
																																			\begin{exampleblock}{Answer: B}\end{exampleblock}
																																			\begin{block}{Explanation}
																																				Data minimisation limits collection to what is necessary. Options A, C, and D violate the principle.
																																			\end{block}
																																		\end{frame}
																																		
																																		\begin{frame}{Q248 --- Model Risk: Responsible AI Principles}
																																			\begin{block}{Question}
																																				Which of the following is NOT one of the widely cited Responsible AI principles?
																																				\begin{enumerate}
																																					\item Fairness.
																																					\item Transparency.
																																					\item Accountability.
																																					\item Maximising short-term shareholder value regardless of harm.
																																				\end{enumerate}  \end{block}
																																				\begin{exampleblock}{Answer: D}\end{exampleblock}
																																				\begin{block}{Explanation}
																																					Responsible AI principles include fairness, transparency, accountability, privacy, and safety. Maximising short-term shareholder value at the expense of harm is contrary to Responsible AI.
																																				\end{block}
																																			\end{frame}
																																			
																																			\begin{frame}{Q249 --- Model Risk: Final Integration}
																																				\begin{block}{Question}
																																					A financial institution is establishing an AI/ML governance framework. Which of the following is a key recommendation from the curriculum?
																																					\begin{enumerate}
																																						\item Build a new framework entirely from scratch.
																																						\item Start from existing model governance frameworks and extend them with AI-specific considerations (data, explainability, ethics, fairness, and GenAI-specific governance).
																																						\item Focus only on technical performance.
																																						\item Ignore regulatory expectations.
																																					\end{enumerate}  \end{block}
																																					\begin{exampleblock}{Answer: B}\end{exampleblock}
																																					\begin{block}{Explanation}
																																						Extending existing governance frameworks with AI-specific elements is recommended. Options A, C, and D are not aligned with best practice.
																																					\end{block}
																																				\end{frame}
																																				
																																																			\end{document}